\section{Introduction}

Machine-learned global weather models now forecast the atmosphere at least as well as physics-based systems on
standard grid-level scores \citep{lam2023graphcast,price2025gencast,lang2024aifs}. ECMWF runs its AI model, AIFS,
operationally, and NOAA runs a GraphCast-based system (AIGFS). In parallel, time-series foundation models such as
TimesFM \citep{das2024timesfm} and Chronos \citep{ansari2024chronos,ansari2025chronos2} forecast any series
zero-shot from its own history.

Many decisions in India depend on a narrower question: what will the thermometer, hygrometer and anemometer at
this airport or this monitoring station read over the next day? Global models answer it only through a grid cell
of roughly 9--28 km, which misses local siting and land use. Foundation models see one sensor at a time and no
weather physics. The classical answer is statistical post-processing of numerical forecasts at stations
\citep{glahn1972mos,gneiting2005emos,rasp2018postprocessing}.

We ask how far a small model can go that is trained only on a station network's own history, optionally with the
global forecasts that exist when the forecast is made. Our model, \pwm, has about one million parameters, more
than 100 times fewer than Chronos-2 (120M) and TimesFM 2.5 (200M). We report its development version by
version under a fixed rule, including the changes that failed.

\paragraph{Contributions.}
\begin{itemize}
\item \textbf{Three north-Indian station leagues} (\S\ref{sec:data}): airports, NOAA's quality-controlled stations
around Delhi, and Delhi's air-quality stations' own weather sensors. The opponents are real archived operational
forecasts of five global systems. Each is scored under two rules: its \emph{latest run}, which is often issued
after our forecast starts, and its \emph{newest run issued before} our forecast starts. Each is also scored raw and
after a trailing local bias correction. TimesFM 2.5, TimesFM 3.0 and Chronos-2 complete the field.
\item \textbf{\pwm} (\S\ref{sec:model}), an inductive station-graph world model, in a pure form (observations only)
and a hybrid form that also reads bias-corrected forecasts from four global models, either as inputs or as a base
it learns to correct.
\item \textbf{A version-by-version record} (\S\ref{sec:versions}) from v0.3 to v0.8. Each change was adopted only if
it beat the default setting on validation by more than two seed standard deviations. Several proposals failed
this test, and we report them alongside the ones that passed.
\item \textbf{Results with losses stated} (\S\ref{sec:results}):
\begin{itemize}
\item at the airports, the hybrid beats AIFS, IFS, ICON, AIGFS and GFS on every variable they forecast under the issued-before
rule, but loses to AIFS's freshest run on temperature and humidity;
\item the pure model beats all three foundation models at the airport and WMO stations;
\item on the CPCB sensors the pure model still trails the foundation models on wind, while the hybrid ties or
beats them;
\item bias-corrected FourCastNet and GenCast Mini, which we ran ourselves, beat the pure model on temperature.
\end{itemize}
\end{itemize}
